% rc-worksheet.tex
\documentclass[11pt, a4paper]{article}

\usepackage[T1]{fontenc}
\usepackage[utf8]{inputenc}
\usepackage{tgpagella}
\usepackage[spanish, es-noshorthands]{babel}
\usepackage{amsmath, amssymb, bm}
\usepackage[margin=2.5cm]{geometry}
\usepackage{fancyhdr}
\usepackage{circuitikz}
\usepackage[most]{tcolorbox}

\pagestyle{fancy}
\fancyhf{}
\setlength{\headheight}{16pt}
\lhead{\textbf{UCB Sede Tarija} | Ing. Mecatrónica}
\rhead{Hoja de Trabajo: Modelado Guiado RC[cite: 4]}
\lfoot{Prof: M.Sc. Ing. Bernardo Quiroga Turdera[cite: 1, 3]}
\renewcommand{\headrulewidth}{0.4pt}
\definecolor{ucbblue}{RGB}{0, 51, 102}

\begin{document}

\begin{tcolorbox}[colback=ucbblue!5!white, colframe=ucbblue, title=\textbf{Artefacto 3: Modelado Guiado RC}]
\end{tcolorbox}

\textbf{Circuito Dado:} Circuito serie simple con fuente $V_s$, resistor $R$ y capacitor $C$, cerrado en $t=0$[cite: 4].
\begin{center}
    \begin{circuitikz}
        \draw (0,0) to[V, l={$V_s$}] (0,2) to[closing switch] (2,2) to[R, l={$R$}] (4,2) to[C, l={$C$}] (4,0) -- (0,0);
    \end{circuitikz}
\end{center}

\vspace{0.5cm}
\textbf{Parte 1 - Ley del Circuito (KVL):}[cite: 4]\\
$V_s - \underline{\hspace{2cm}} - \underline{\hspace{2cm}} = 0$[cite: 4]

\vspace{0.5cm}
\textbf{Parte 2 - Ley del Resistor (Ohm):}[cite: 4]\\
$v_R = \underline{\hspace{4cm}}$[cite: 4]

\vspace{0.5cm}
\textbf{Parte 3 - Ley del Capacitor:}[cite: 4]\\
$i = \underline{\hspace{4cm}}$[cite: 4]

\vspace{0.5cm}
\textbf{Parte 4 - Sustitución:}[cite: 4]\\
$R \left( \underline{\hspace{4cm}} \right) + v_C = V_s$[cite: 4]

\vspace{0.5cm}
\textbf{Parte 5 - Modelo Final:}[cite: 4]\\
$\underline{\hspace{2cm}} \frac{dv_C}{dt} + v_C = \underline{\hspace{2cm}}$[cite: 4]

\vspace{0.5cm}
\textbf{Parte 6 - Explicación Conceptual:}[cite: 4]\\
El término $\frac{dv_C}{dt}$ representa \underline{\hspace{8cm}}.[cite: 4] \\
La ecuación es de primer orden porque \underline{\hspace{8cm}}.[cite: 4] \\
La ecuación describe cómo \underline{\hspace{4cm}} evoluciona debido a \underline{\hspace{4cm}}.[cite: 4]

\vspace{0.5cm}
\textbf{Parte 7 - Condiciones Iniciales y Finales:}[cite: 4]\\
Si $v_C(0^-) = \underline{\hspace{2cm}}$, entonces por continuidad $v_C(0^+) = \underline{\hspace{2cm}}$.[cite: 4]\\
En estado estacionario final de DC, $\frac{dv_C}{dt} = \underline{\hspace{2cm}}$. Por ende inferimos que $v_C(\infty) = \underline{\hspace{2cm}}$.[cite: 4]

\end{document}
